% Canonical entry point for the current ACM-formatted manuscript.
% Edit h2snn_manuscript.tex; material below \endinput is an inactive legacy draft.
\input{h2snn_manuscript.tex}
\endinput
\documentclass{article}


\usepackage[preprint]{neurips_2024}


\usepackage[utf8]{inputenc}
\usepackage[T1]{fontenc}
\usepackage{hyperref}
\usepackage{url}
\usepackage{booktabs}
\usepackage{amsfonts}
\usepackage{amsmath}
\usepackage{amssymb}
\usepackage{nicefrac}
\usepackage{microtype}
\usepackage{xcolor}
\usepackage{graphicx}
\usepackage{algorithm}
\usepackage{algpseudocode}
\usepackage{multirow}
\usepackage{array}
\usepackage{makecell}

\newcommand{\eqdef}{\coloneqq}
\newcommand{\R}{\mathbb{R}}

\title{H2S-NN: Heterogeneous Sheaf-Hodge Neural Networks \\ with Learned Schema-to-Cell Discovery}

% Fill in real author / affiliation details before submission.
\author{
  Peeradon Sarnkaew \\
  \texttt{peeradonsarnkaew46@gmail.com} \\
  % Institution / department placeholder --- update before submission.
}

\begin{document}

\maketitle

\begin{abstract}
Heterogeneous information networks --- graphs with multiple node and edge types --- are frequently
also heterophilic: nodes connected by an edge need not share a label or even a similar semantic
role. Ordinary graph neural networks handle this combination poorly, because standard message
passing lives on the graph's 1-skeleton and either over-smooths dissimilar neighbors or ignores
higher-order typed relational patterns altogether. Two existing lines of work each solve half of
this problem. Heterogeneous sheaf neural networks (notably HetSheaf, \citealp{braithwaite2024hetsheaf})
learn typed restriction maps that let connected nodes disagree in raw feature space while remaining
compatible after a learned, relation-specific transport --- but they operate only on the fixed graph
structure and never construct higher-order topology. Meta-path search methods
(LMSPS~\citep{yan2024lmsps}, PMMM~\citep{jiang2022pmmm}, AMP-HGNAS, MCTS-HGNN, and others) automatically
discover which typed paths are informative, but every method in this family treats a selected path
purely as an aggregation channel: the path is followed to collect a neighborhood and then discarded.
We introduce \textbf{H2S-NN} (Heterophilic Heterogeneous Sheaf-Hodge Neural Network), which unifies
these two directions. Concrete meta-path \emph{instances} --- not merely their type-level schemas
--- are lifted into two-dimensional cells of a typed cellular complex, and a learned sheaf on that
complex (restriction maps on nodes, edges, and faces) drives Hodge-style diffusion over both node and
edge cochains through the coboundary operators $\delta_0,\delta_1$ and their Laplacians $L_0,L_1$,
in the spirit of cellular-sheaf diffusion for homogeneous graphs~\citep{bodnar2022neural}. We further
introduce a schema-to-cell discovery extension that removes the dependence on hand-specified
meta-paths: a training-free, topology-aware score identifies a compact pool of candidate path
schemas, and a hierarchical pair of differentiable gates --- a schema-level gate $\alpha$ and an
instance-level gate $\beta$ --- learns which schemas and which concrete instances should shape the
diffusion operator, stabilized by a partial-schema-propagation mechanism. On node classification
(ACM, DBLP, IMDB), graph classification (MUTAG, PROTEINS, ENZYMES, NCI1), and heterogeneous link
prediction (LastFM, PubMed, Amazon, YouTube), H2S-NN reports the strongest numbers in our assembled
comparison against homogeneous, heterogeneous, and sheaf-based baselines (Section~\ref{sec:results});
we are explicit throughout about which parts of this evidence base --- repeated-seed statistics, the
full ablation matrix, and the discovery extension's own evaluation --- still need to be completed
before the claims below should be read as fully validated (Sections~\ref{sec:exp-setup}
and~\ref{sec:discussion}).
\end{abstract}

\section{Introduction}
\label{sec:intro}

\subsection{Motivation}

We start from a heterogeneous information network: multiple node types, multiple relation types,
and --- in the settings that motivate this work --- often strong heterophily, meaning connected
nodes may not share labels or semantics. This combination creates two distinct problems for
ordinary GNNs. First, simple neighborhood aggregation can over-smooth dissimilar nodes, since it
implicitly assumes that connected nodes should become more similar after message passing. Second,
ordinary message passing lives mainly on the graph's 1-skeleton --- its nodes and edges --- so it
does not explicitly model higher-order typed relational patterns, even when such patterns (a paper
written by a specific pair of co-authors, a specific citation chain) are exactly the structure a
domain expert would point to as meaningful. Most heterogeneous GNN methods address the second
problem by introducing meta-paths as additional graph views or aggregation channels
\citep{wang2019han,fu2020magnn}: a meta-path such as \emph{paper--author--paper} defines a new,
derived adjacency that a model can attend over. This is a real improvement, but it treats the
meta-path as nothing more than a routing rule.

\subsection{From semantic view to topological cell}

H2S-NN takes a different position. Rather than using a meta-path as a semantic aggregation
channel, we use each concrete instance of a meta-path as a higher-order \emph{topological} object.
Concretely, a path instance such as $a_1 \to p_7 \to a_2$ becomes a face of a cellular complex;
faces participate in the coboundary operator $\delta_1$; $\delta_1$ contributes to the up-Laplacian
term of $L_1$; and $L_1$ directly changes the Hodge diffusion geometry over edge signals. The shift
in one sentence: \emph{meta-path as semantic view} becomes \emph{meta-path instance as topological
cell}. This is not a relabeling of existing search methods --- it changes what the selected
structure is used for. Every existing automatic meta-path/schema search method we are aware of
(LMSPS~\citep{yan2024lmsps}, PMMM~\citep{jiang2022pmmm}, AMP-HGNAS, MCTS-HGNN, and several others)
uses a selected path the same way HAN~\citep{wang2019han} does: to route a neighborhood into an
attention-weighted sum. None of them make the path itself, as opposed to the neighbors at the end
of it, a meaningful structural object that reshapes the model's operator.

\subsection{A precise statement of what is (and is not) new}

An earlier draft of this project claimed novelty from ``working without predefined meta-paths.'' We
verified this claim against the literature and found it to be false: LMSPS, AMP-HGNAS, MCTS-HGNN,
and PMMM all already discover schemas automatically, without a human-specified meta-path set. We
therefore do \emph{not} lead with this framing, and we flag its removal explicitly here rather than
let a reviewer discover the overclaim independently. The real, verified gap is narrower and more
specific: it is in \emph{what the selected structure becomes} --- an aggregation channel for every
prior method, versus a two-dimensional cell that reshapes $\delta_1$ and $L_1$ for H2S-NN --- and in
\emph{how it acts on the model} --- a weighted sum over neighbors, versus a term inside a
matrix-free Hodge Laplacian. This distinction is also why a direct comparison against HetSheaf ---
the closest sheaf-based baseline --- is not entirely fair to HetSheaf: HetSheaf operates on the
observed 1-skeleton only, with no notion of a face at all.

\subsection{Contributions}

\begin{enumerate}
  \item \textbf{Operator-level higher-order modeling.} We show how to turn selected meta-path
  instances into two-dimensional cells of a typed cellular complex such that a learned edge-to-face
  sheaf directly reshapes the up-Laplacian term of $L_1$, rather than merely providing an
  additional aggregation channel.
  \item \textbf{Full Hodge diffusion on a heterogeneous complex.} We define matrix-free coboundary
  operators $\delta_0,\delta_1$ and Laplacians $L_0,L_1$ over both node and edge cochains, and run
  $K$-step Hodge diffusion with learnable step sizes.
  \item \textbf{Hierarchical, jointly learned schema and instance gating.} A schema-level gate
  $\alpha_\sigma$ and an instance-level gate $\beta_f$ are trained together and multiplied into the
  weighted up-Laplacian, giving both a coarse (which meta-path template matters) and a fine (which
  specific instance of it matters) control signal.
  \item \textbf{Training-free, topology-aware schema scoring.} Candidate schema quality is
  estimated from an endpoint-label-discrepancy statistic computed via label propagation over
  training labels only, avoiding the cost of a full model forward pass per candidate schema and,
  by construction, any use of validation or test labels during schema selection.
\end{enumerate}

\subsection{Scope of this paper and how its two contributions are evidenced}

This paper presents H2S-NN as a \emph{unified} contribution spanning both the base sheaf-Hodge
model (contributions 1--2 above) and the schema-to-cell discovery extension (contributions 3--4).
We take this risk deliberately, but we manage it explicitly rather than let the two halves blur
into one inflated claim: the headline tables in Section~\ref{sec:results} report the \emph{full,
discovery-enabled} configuration (Section~\ref{sec:ablation-schema}'s clarification), and
Section~\ref{sec:discussion} returns to precisely what is and is not yet evidenced about that
configuration before recommending how the two contributions should be weighted in a submission.

\section{Related Work}
\label{sec:related}

\subsection{Heterogeneous GNNs without meta-path search}

HAN~\citep{wang2019han} and MAGNN~\citep{fu2020magnn} use expert-defined meta-paths as semantic
message-passing views, attending over a fixed, hand-designed set of relational patterns.
HGT~\citep{hu2020hgt} instead uses typed node- and edge-level attention without any hand-designed
meta-path, learning heterogeneity directly through typed transformer weights. Neither family turns
a concrete path instance into an operator-shaping higher-order cell; both remain message-passing
models over the graph's 1-skeleton.

\subsection{Meta-path and schema search methods}

A substantial and active literature automates the choice of which meta-paths or schemas to use:
LMSPS~\citep{yan2024lmsps} and AMP-HGNAS learn differentiable or NAS-style selection weights over
precomputed path features; MCTS-HGNN applies Monte Carlo tree search directly over the schema
space; PMMM~\citep{jiang2022pmmm} and the related Meta-Multigraph Search search meta-multigraphs
with a partial message propagation technique for stabilizing the resulting bilevel optimization;
HG-Search and SchemaWalk apply, respectively, a multi-stage and an inductive
reinforcement-learning-flavored search; RMS-HRec, SMCD, HeteroSample, EvoPath, TAEG-HG, FHGE,
HISE-KT, ESim, MP-GNN, and LPSMS round out the broader family with recommendation-, LLM-, and
sampling-oriented variants. The unifying limitation is that every one of these methods ---
regardless of how sophisticated its search procedure --- uses a selected meta-path as an
\textbf{aggregation channel}: the path is followed to collect a neighborhood and then discarded; the
higher-order relational pattern is thrown away after serving as a routing rule.

\subsection{Sheaf neural networks}

Neural Sheaf Diffusion~\citep{bodnar2022neural} established the core connection between learned
cellular-sheaf restriction maps, heterophily, and the avoidance of naive oversmoothing, but for
homogeneous graphs with no notion of node or edge types. HetSheaf~\citep{braithwaite2024hetsheaf}
extends this to the heterogeneous setting, learning typed sheaf restriction maps and handling
heterophily through per-relation local spaces, and further introduces SheafPool, a universal
stalk-space graph-level readout invariant to local changes of basis; it is, by a wide margin, the
closest existing baseline to H2S-NN and we treat it as such throughout this paper (its own
evaluation is built on the Heterogeneous Graph Benchmark, HGB~\citep{lv2021hgb}, which we also
draw on for baseline numbers and evaluation conventions). Other sheaf applications (DeepSN, SHN)
apply sheaf-theoretic machinery to different problems again without constructing higher-order
topology. Across the sheaf literature we are aware of, H2S-NN is, to our knowledge, the only method
that uses both the up-Laplacian $L_1$ and the coboundary operator $\delta_1$; HetSheaf and the
other sheaf methods operate with typed restriction maps on the observed 1-skeleton, with no face
topology, no coboundary operator beyond $\delta_0$, and no Hodge decomposition on edge cochains.

\subsection{Heterophily-aware heterogeneous models}

A separate, more recent line of work addresses heterophily in heterogeneous graphs directly:
HDHGR~\citep{guo2023hdhgr} proposes a homophily-oriented heterogeneous graph rewiring approach
guided by a meta-path-induced homophily metric; H2SGNN~\citep{lu2024h2sgnn} applies local
independent filtering and global hybrid filtering from a spectral perspective to address
heterogeneity and heterophily jointly; LatGRL~\citep{shen2024latgrl} defines meta-path-level and
node-level semantic homophily ratios and couples global structure/attribute similarity mining with
an adaptive dual-frequency fusion mechanism; HGMS~\citep{wang2025hgms} is a homophily-aware
heterogeneous graph contrastive learning framework that augments views via connection-strength-guided
edge dropping; and the Cross-Type Homophily line, CTHGE~\citep{tao2025cthge}, proposes an explicit
Cross-Type Homophily Ratio and a corresponding graph-editing method. These methods address
heterophily through structural, spectral, or contrastive mechanisms, not through weighted typed
path instances acting as Hodge faces, and they form part of our baseline set
(Section~\ref{sec:exp-setup}) rather than a competing higher-order-topology approach.

\subsection{Positioning summary}

\begin{table}[h]
  \centering
  \caption{H2S-NN's position relative to five adjacent research lines.}
  \label{tab:positioning}
  \small
  \begin{tabular}{p{3.1cm}p{4.6cm}p{4.6cm}}
    \toprule
    \textbf{Research line} & \textbf{What it solves} & \textbf{Gap H2S-NN targets} \\
    \midrule
    HAN / MAGNN style meta-path HGNNs & Use expert-defined meta-paths as semantic message-passing views. & They do not turn each concrete path instance into an operator-shaping higher-order cell. \\
    HGT style transformer HGNNs & Use typed node and edge attention without hand-designed meta-paths. & They remain edge/message-passing models rather than Hodge models over faces. \\
    Meta-path search methods (LMSPS, PMMM, \ldots) & Search useful semantic paths or meta-multigraphs. & Selected schemas usually become aggregation channels, not cells in a learned Hodge operator. \\
    HetSheaf & Learns heterogeneous sheaf maps and handles heterophily through local spaces. & Closest baseline, but H2S-NN aims to add path-induced higher-order structure and $\delta_1$. \\
    HDHGR / H2SGNN / LatGRL / HGMS / CTHGE & Recent heterophily-aware heterogeneous models with strong graph-level mechanisms. & They address heterophily, but not through weighted typed path instances as Hodge faces. \\
    \bottomrule
  \end{tabular}
\end{table}

\section{Method}
\label{sec:method}

\subsection{Preliminaries and notation}

Let $N$ and $M$ denote the number of nodes and (non-closure) edges in the heterogeneous graph, $d$
the shared stalk dimension used for node, edge, and face states, $d_r$ the relation-embedding
dimension, and $d_h$ the hidden dimension of the sequence encoder used for edge-to-face maps. Each
node $v$ has a type $\phi(v)$; each relation $r$ belongs to a relation set $R$ and has a learned
embedding $z^{rel}_r \in \R^{d_r}$. A schema $\sigma$ is a typed path template of length $\ell$
(number of edges); a face $f$ is one concrete instance of a schema, and $F_s$ denotes the number of
face instances materialized for schema $s$. We write $H^{(0)}$ (equivalently $h^{(0)}$) for the
initial node encoding, $F^{src}_e, F^{tgt}_e \in \R^{d\times d}$ for the source/target node-to-edge
restriction maps of edge $e$, $s^{(0)}$ for the initial edge cochain, $F_{e\triangleright f}$ for
the edge-to-face restriction map, $\delta_0,\delta_1$ for the coboundary operators and $L_0, L_1$
for their associated Hodge Laplacians, $K$ for the number of Hodge-diffusion (Euler) steps with
step sizes $\eta_0,\eta_1$, $\alpha_\sigma$ and $\beta_f$ for the schema- and instance-level
discovery gates, and $x^*, s^*$ for the post-diffusion node and edge states used at readout.

\subsection{Typed node encoding}

Node features are encoded into the shared stalk dimension by a type-specific encoder:
\begin{equation}
  h^{(0)}_v = \mathrm{GELU}\big(\mathrm{LayerNorm}(\mathrm{Linear}_{\phi(v)}(x_v)) + \mathrm{TypeEmbed}[\phi(v)]\big) \in \R^d.
  \label{eq:node-enc}
\end{equation}
Each node type owns an independent two-layer encoder, and a single shared, learnable type embedding
is added to every node of that type. For graphs that ship without node features at all --- LastFM
and AmazonBook, in our link-prediction experiments --- a plain linear encoder of the form
$\mathrm{Linear}(0, d)$ would collapse every node of a type to an identical vector (its bias term);
we instead fall back to a free, per-node embedding table $\mathrm{Embedding}(N, d)$, so that
featureless nodes still receive individually learnable initial states.

\subsection{Node-to-edge sheaf restriction maps}

For every directed, relation-typed edge $e = (u \to_r v)$, a pair of restriction maps
$F^{src}_e, F^{tgt}_e$ is produced by relation-specific two-layer GELU MLPs taking
$[h^{(0)}_u \Vert h^{(0)}_v \Vert z^{rel}_r]$ as input. By default these maps are constrained to be
\emph{scaled orthogonal} rather than free $d\times d$ matrices: an unconstrained parameterization is
mapped to an orthogonal matrix by one of four supported constructions --- matrix exponential,
Cayley transform, a product of Householder reflectors (the default), or, for $d\in\{2,3\}$,
closed-form Euler rotations --- and then multiplied by a learned scalar scale,
$\exp(\mathrm{clamp}(\text{raw}, \max=2.0))$. We are careful to describe these as \emph{scaled
orthogonal} restriction maps throughout the paper rather than simply ``orthogonal,'' since the
scaling breaks strict orthogonality. These maps define the first coboundary operator on the initial
node encoding,
\begin{equation}
  (\delta_0 h^{(0)})_e = F^{tgt}_e h^{(0)}_{tgt(e)} - F^{src}_e h^{(0)}_{src(e)},
  \label{eq:delta0-h0}
\end{equation}
which is computed once and reused as an input feature to the edge initializer below --- never as
the diffusion's initial edge state directly. An ablation variant, \texttt{ablate\_identity\_maps},
replaces both restriction maps with the identity, degenerating $\delta_0$ into a plain, untyped
node-difference operator; we use this as one of our minimum ablations
(Section~\ref{sec:ablations}).

\subsection{Edge cochain initialization}

The initial edge cochain is produced by a shared two-layer MLP:
\begin{equation}
  s^{(0)}_e = \mathrm{MLP}_{init}\big([h^{(0)}_u \Vert h^{(0)}_v \Vert z^{rel}_r \Vert (\delta_0 h^{(0)})_e]\big) \in \R^d.
\end{equation}
This is a deliberate design choice, not an oversight: setting $s^{(0)} = \delta_0 H^{(0)}$ directly
would zero out exactly the harmonic component of the initial edge signal that Hodge diffusion is
meant to preserve and evolve. $\delta_0 H^{(0)}$ is therefore only ever consumed as one input
feature among several, controlled by the \texttt{use\_delta0\_init} flag (default on).

\subsection{Edge-to-face maps via sequence encoding}

Given a schema's set of concrete path instances, each of length $\ell$ edges, we encode the ordered
sequence of boundary edges with a single-layer LSTM plus an attention head. For a chunk of $B$
faces (chunked at \texttt{face\_chunk\_size}, default 4096, to bound peak memory), token features
$\mathrm{token}_i = \mathrm{MLP}_g([h^{(0)}_{src(e_i)} \Vert h^{(0)}_{tgt(e_i)} \Vert
z^{rel}_{r_i}])$ are passed through the LSTM to produce hidden states $h_1,\dots,h_\ell$; a learned
attention vector $a$ produces position weights $\alpha_i^{pos} = \mathrm{softmax}(h_i \cdot a /
\sqrt{d_h})$ over the path position (superscripted to disambiguate from the schema-level discovery
gate of the same Greek letter). The resulting edge-to-face restriction map is
\begin{equation}
  F_{e_i\triangleright f} = \alpha_i^{pos} \cdot \Psi(h_i), \qquad \Psi(h_i) = \mathrm{reshape}\big(\mathrm{Linear}(h_i),\, [d,d]\big).
  \label{eq:edge-to-face}
\end{equation}
We flag this formula precisely because it differs from an earlier draft of the internal
mathematical specification, which stated that the attention weights are used only in the
coboundary regularizer and \emph{not} in the map construction itself. The implementation multiplies
attention directly into the map, and we describe the model as it is actually implemented. These
maps are computed once per forward pass and cached, then reused across all $K$ diffusion steps, so
$L_0$ and $L_1$ are fixed linear operators for the duration of the Euler rollout described in
Section~\ref{sec:diffusion}.

\subsection{Closure edges as a closed-form correction}

For each face, the underlying cellular-complex construction requires closing the open path instance
with a synthetic \emph{closure edge} from its final node back to its initial node, so that the
face's boundary forms a genuine closed loop. Rather than materializing this as an additional 1-cell
with its own learned cochain state, we implement it as a closed-form correction added into $L_0$
via index-add over the closure endpoints:
\begin{equation}
  \Delta_{L_0} = \texttt{closure\_scale}^2 \cdot (x_{tgt(f)} - x_{src(f)}),
\end{equation}
where the squared dependence on \texttt{closure\_scale} follows from $L_0 = \delta_0^\top \delta_0$
being quadratic in the restriction map, here a frozen, uniformly scaled identity. With the default
\texttt{closure\_scale}$=10^{-6}$, the effective coupling introduced by closure edges is
$10^{-12}$ --- functionally negligible, but present and load-bearing for the closed-boundary claim,
and consciously distinct from allowing closure edges to affect $\delta_1$ at all, which they do
not, by design.

\subsection{Coboundary operators and Hodge Laplacians}
\label{sec:diffusion}

All coboundary and Laplacian operators are implemented matrix-free --- never materialized as dense
matrices --- using batched per-edge and per-face matrix--vector products and scatter-adds:
\begin{align}
  \delta_0(x)_e &= F^{tgt}_e x_{tgt(e)} - F^{src}_e x_{src(e)} \\
  L_0(x) &= \delta_0^\top \delta_0\, x + \Delta_{L_0}(x) \\
  \delta_1(s)_f &= \sum_{e \in \partial f} \eta(e,f)\, F_{e\triangleright f}\, s_e, \qquad \eta(e,f)\in\{+1,-1\} \\
  L_1(s) &= \delta_0 \delta_0^\top s + \delta_1^\top \delta_1\, s
\end{align}
The orientation sign $\eta(e,f)$ is implemented as a general, explicit, dense-reference-tested
tensor, but in practice it is always $+1$ in this project's data, since our path-instance
enumeration always walks forward through each relation's stored source-to-destination direction. We
nonetheless treat the general sign machinery as load bearing, since future extensions using
undirected or backward-traversed boundary edges would require it. When schema/instance discovery
weights are active, the up-Laplacian term generalizes to
\begin{equation}
  L_1^{up}(s) = \sum_{\sigma} \alpha_\sigma \sum_{f \in \sigma} \beta_f\, (\delta_1^f)^\top \delta_1^f\, s,
  \label{eq:l1up}
\end{equation}
which is the operator through which the learned discovery gates of Section~\ref{sec:discovery} act
directly on diffusion geometry --- the property that distinguishes H2S-NN's use of discovered
structure from every prior meta-path search method (Section~\ref{sec:related}).

\subsection{Hodge diffusion}

Node and edge states are updated jointly for $K$ steps of explicit Euler integration with
learnable step sizes:
\begin{equation}
  x \leftarrow x - \eta_0 L_0(x), \qquad s \leftarrow s - \eta_1 L_1(s), \qquad k = 1,\dots,K,
\end{equation}
where $\eta_0,\eta_1$ are reparameterized through a softplus so that they remain positive while
still being trained end to end. The resulting states $x^*,s^*$ are each passed through a
\texttt{LayerNorm} before being handed to readout.

\subsection{Readout heads}

The base node-classification readout forms, for each target-type node $v$, the concatenation
\begin{equation}
  z_v = \big[x^*_v \,\Vert\, h^{(0)}_v \,\Vert\, \mathrm{pool\_edges\_to\_nodes}(x^*,s^*)_v\big]
\end{equation}
of the final node state, the initial node state, and a mean pooling of incident edge states over
both endpoints, and passes $z_v$ through a small MLP classifier. Because this is a unified paper
also covering link prediction and graph classification, two further readout heads share the same
encoder/diffusion backbone: a link-prediction readout with a swappable dot-product, DistMult, or
concatenation-plus-MLP decoder over node-pair states; and a graph-level pooling module with several
swappable heads --- naive sum-, mean-, and max-pooling baselines, an $O(d)$-invariant \textsc{GramPool}
(which exploits that restriction maps are always orthogonal in every configuration we run), a
gauge-fixing attention-barycenter \textsc{SheafPool} baseline, and a $GL(d)$-invariant \textsc{S2Pool}.
We treat the gauge-invariance argument behind \textsc{GramPool}/\textsc{S2Pool} as a secondary
methodological contribution: because sheaf restriction maps are only well-defined up to a choice of
local frame (``gauge'') at each stalk, a graph-level readout that is not itself invariant to that
choice risks encoding an artifact of the parameterization rather than a genuine structural signal,
and \textsc{S2Pool} is constructed to be invariant to general invertible gauge transformations
rather than only orthogonal ones.

\subsection{Training objective}

The full training loss combines a supervised classification term with four sheaf-geometric
regularizers and, once discovery gates are active, six further discovery-specific terms:
\begin{equation}
  \mathcal{L} = \mathcal{L}_{cls} + \mathcal{L}_{smooth0} + \mathcal{L}_{smooth1} + \mathcal{L}_{consist} + \mathcal{L}_{coboundary} + \mathcal{L}_{discovery}.
\end{equation}

\begin{table}[h]
  \centering
  \caption{Sheaf-geometric regularizer terms and their default weights.}
  \label{tab:loss-terms}
  \small
  \resizebox{\textwidth}{!}{%
  \begin{tabular}{lll}
    \toprule
    \textbf{Term} & \textbf{Formula} & \textbf{Default $\lambda$} \\
    \midrule
    Classification & Cross-entropy (single-label) or sigmoid + BCE (multi-label) on the training mask & --- \\
    Smooth0 & $\lambda_0 \cdot \mathrm{mean}(\lVert\delta_0 x^*\rVert^2)$ & $0.0$ \\
    Smooth1 & $\lambda_1 \cdot \mathrm{mean}(\lVert\delta_1 s^*\rVert^2)$ & $10^{-3}$ \\
    Consistency & $\lambda_c \cdot \mathrm{mean}(\lVert s^* - \delta_0 x^*\rVert^2)$ & $10^{-2}$ \\
    Coboundary & $\lambda_{co} \cdot \mathrm{mean}(\lVert\delta_1 \delta_0 H^{(0)}\rVert^2)$ & $10^{-2}$ \\
    \bottomrule
  \end{tabular}%
  }
\end{table}

The coboundary term is deliberately evaluated at the \emph{initial} encoding $H^{(0)}$ rather than
the diffused state $x^*$: this softly penalizes violation of $\delta_1 \circ \delta_0 = 0$ without
competing, through the same tensor, against the diffusion objective's own gradient signal.

Once schema/instance discovery is enabled and past its warm-up period, six further terms are added
--- schema sparsity ($\lambda=10^{-3}$), instance sparsity ($10^{-4}$), redundancy ($10^{-3}$),
coverage ($10^{-3}$), cost ($10^{-4}$), and gate entropy ($10^{-4}$). One implementation-precision
point worth stating: the instance-sparsity term must read a gradient-attached copy of the mean
instance weights, not the detached copy used purely for logging, so that the sparsity regularizer
actually trains the instance gate's parameters.

\textbf{Multi-label handling.} IMDB's HGB preprocessing assigns each movie a multi-hot genre label
rather than a single class. The data loader detects this via a dimensionality check on the target
tensor and threads an \texttt{is\_multilabel} flag into the trainer, which switches to a sigmoid
output and binary cross-entropy loss with multi-label F1 metrics instead of the single-label path
--- matching HGB's own stated convention for multi-label datasets~\citep{lv2021hgb}.

\subsection{Schema-to-cell discovery}
\label{sec:discovery}

The auto-discovery pipeline proceeds in five stages.

\textbf{(A) Schema graph.} A type-level adjacency is built from the set of relation tuples, giving,
for each node type, the set of relations that can be followed to each neighboring type.

\textbf{(B) Candidate generation.} Typed walks are enumerated by depth-first search over the schema
graph, bounded by \texttt{min\_schema\_length}/\texttt{max\_schema\_length} (2/4 by default),
optionally required to both start and end at the target node type, and capped at
\texttt{max\_schema\_candidates} (128 by default) to keep the candidate pool tractable.

\textbf{(C) Topology-aware, training-free scoring.} Rather than evaluate every candidate schema
with a full model forward pass, each candidate is scored from purely structural and
label-propagation statistics: an \emph{endpoint label discrepancy} computed by propagating
\textbf{training labels only} to estimate how dissimilar the predicted labels at a path's two
endpoints tend to be --- schemas connecting label-dissimilar endpoints are exactly the
heterophilic, informative patterns we want to surface --- alongside coverage, relation and type
diversity, and average branching factor.

\textbf{(D) Greedy topology-aware pool selection.} Schemas are added one at a time by a greedy
forward-selection rule,
\begin{equation}
  \mathrm{score}(\sigma \mid \text{selected}) = \frac{\mathrm{disc}(\sigma)^{\alpha_1}\, \mathrm{cov}(\sigma)^{\alpha_2}\, \mathrm{div}(\sigma)^{\alpha_3}}{1 + \max_{\sigma'\in\text{selected}}\mathrm{redundancy}(\sigma,\sigma')},
\end{equation}
with redundancy measured as the Jaccard overlap between sampled endpoint sets, until
\texttt{selected\_schema\_k} (8 by default) schemas are retained.

\textbf{(E) Face materialization.} For each selected schema, concrete path instances are sampled,
closure-edge pairs are attached, and a per-schema structural feature vector is computed and
consumed directly by the discovery losses above.

\textbf{Hierarchical gates.} The selected structure is weighted by a product of a schema-level and
an instance-level gate:
\begin{align}
  \alpha_\sigma &= \mathrm{sigmoid}(\mathrm{logit}_\sigma), \quad \text{warm-started from } \log(\text{topology score}_\sigma), \\
  \beta_f &= \mathrm{sigmoid}\big(\mathrm{MLP}([h^{(0)}_{v_0} \Vert h^{(0)}_{v_\ell} \Vert \text{schema\_embedding}_\sigma])\big);
\end{align}
the effective per-face weight used in $L_1^{up}$ (Eq.~\ref{eq:l1up}) is $\alpha_\sigma \cdot
\beta_f$. The warm start means schemas the structural stage already flagged as high-discrepancy
begin training with high $\alpha$ and must be \emph{proven wrong} by the discovery losses and
downstream task loss, rather than discovered entirely from a cold start.

\textbf{Partial schema propagation.} During training, each schema's gate is independently
Bernoulli-masked and unbiasedly rescaled, $\alpha \leftarrow \alpha \cdot \mathrm{Bernoulli}(p)/p$
with \texttt{partial\_schema\_probability} $p=0.5$ by default, adapted from PMMM's partial message
propagation to address the coupled bilevel-optimization instability that arises when
architecture-like weights and the rest of the model's parameters are updated in the same gradient
step.

Algorithm~\ref{alg:discovery} summarizes the full discovery pipeline.

\begin{algorithm}[t]
  \caption{Schema-to-cell discovery (training-time forward pass)}
  \label{alg:discovery}
  \begin{algorithmic}[1]
    \Require Type graph, relation tuples, training labels $Y_{train}$, candidate budget, pool size $k$
    \State Build schema graph from relation tuples \Comment{Stage A}
    \State Enumerate candidate schemas via bounded DFS \Comment{Stage B}
    \For{each candidate schema $\sigma$}
      \State Sample path instances; estimate endpoint label discrepancy from propagated $Y_{train}$ \Comment{Stage C}
    \EndFor
    \State Greedily select top-$k$ schemas by topology-aware score (Eq.~9) \Comment{Stage D}
    \State Materialize face groups for the selected pool; attach closure edges \Comment{Stage E}
    \State Warm-start $\alpha_\sigma \gets \log(\text{topology score}_\sigma)$
    \While{training}
      \State Compute $\beta_f$ from current node/schema embeddings
      \State Apply partial schema propagation to $\alpha$
      \State Assemble $L_1^{up}$ (Eq.~\ref{eq:l1up}) and run Hodge diffusion
      \State Backpropagate task loss + sheaf regularizers + discovery losses
    \EndWhile
  \end{algorithmic}
\end{algorithm}

\subsection{Implementation notes for methodological precision}

We collect here every point at which the implementation's actual behavior differs from an earlier
version of the internal mathematical specification, so that the rest of this paper describes the
model exactly as built. Attention \emph{is} folded multiplicatively into the edge-to-face
restriction maps (Eq.~\ref{eq:edge-to-face}). Restriction maps are \emph{scaled orthogonal}, not
strictly orthogonal. Closure edges carry no relation embedding of their own. A 6-dimensional
per-instance structural feature vector is computed during face materialization but is not consumed
by the instance gate or any loss term we could locate. A \texttt{hard\_concrete} gate variant is a
listed configuration option but is not yet implemented; we do not report results for it.

\section{Experimental Setup}
\label{sec:exp-setup}

\subsection{Datasets and tasks}

\begin{table}[h]
  \centering
  \caption{Datasets, tasks, and the exact loader verified against the source for each.}
  \label{tab:datasets}
  \small
  \resizebox{\textwidth}{!}{%
  \begin{tabular}{lp{5.3cm}p{6.4cm}}
    \toprule
    \textbf{Task} & \textbf{Datasets} & \textbf{Loader (verified in code)} \\
    \midrule
    Node classification & ACM (target: paper), DBLP (target: author), IMDB (target: movie); manual meta-path schemas per dataset & \texttt{torch\_geometric.datasets.HGBDataset} --- the official HGB preprocessing, not the older HAN-paper preprocessing of the same-named datasets \\
    Link prediction & LastFM, AmazonBook, PubMed, YouTube & \texttt{torch\_geometric.datasets.LastFM}/\texttt{.AmazonBook} for the first two; \textbf{no loader exists for PubMed or YouTube at all} (Sec.~\ref{sec:gaps}) \\
    Graph classification & MUTAG (stratified 10-fold), PROTEINS, ENZYMES, NCI1 (precomputed 10-fold) & \texttt{torch\_geometric.datasets.TUDataset} --- standard preprocessing, no ambiguity \\
    Node clustering & ACM, Freebase, AMiner, IMDB & Not yet implemented (Sec.~\ref{sec:gaps}) \\
    \bottomrule
  \end{tabular}%
  }
\end{table}

We verified the exact loader for every task directly against the source rather than assume
same-named datasets mean the same preprocessing as in other papers --- HGB's own motivation section
makes exactly this point about incompatible numbers across papers using differently-preprocessed,
same-named heterogeneous benchmarks~\citep{lv2021hgb}. Table~\ref{tab:dataset-stats-node} and
Table~\ref{tab:dataset-stats-graph} give the resulting statistics; Table~\ref{tab:dataset-stats-link}
reports a real consequence of this check for our link-prediction results (Sec.~\ref{sec:gaps}).

\begin{table}[h]
  \centering
  \caption{Node classification dataset statistics (HGB's official Table 2~\citep{lv2021hgb} --- matches our \texttt{HGBDataset} loader exactly).}
  \label{tab:dataset-stats-node}
  \small
  \begin{tabular}{lcccccc}
    \toprule
    \textbf{Dataset} & \#Nodes & \#Node types & \#Edges & \#Edge types & Target & \#Classes \\
    \midrule
    DBLP & 26{,}128 & 4 & 239{,}566 & 6 & author & 4 \\
    IMDB & 21{,}420 & 4 & 86{,}642 & 6 & movie & 5 (multi-label) \\
    ACM & 10{,}942 & 4 & 547{,}872 & 8 & paper & 3 \\
    Freebase & 180{,}098 & 8 & 1{,}057{,}688 & 36 & book & 7 \\
    \bottomrule
  \end{tabular}
\end{table}
Freebase is listed for completeness since HGB includes it among its node-classification datasets,
but we only use it for the (currently unimplemented) node-clustering comparison of
Table~\ref{tab:clustering}, not for node classification.

\begin{table}[h]
  \centering
  \caption{Graph classification dataset statistics (standard TUDataset preprocessing).}
  \label{tab:dataset-stats-graph}
  \small
  \begin{tabular}{lccccl}
    \toprule
    \textbf{Dataset} & \#Graphs & \#Classes & Avg.\ \#Nodes & Avg.\ \#Edges & Domain \\
    \midrule
    MUTAG & 188 & 2 & 17.9 & 19.8 & Molecules (mutagenicity) \\
    PROTEINS & 1{,}113 & 2 & 39.1 & 72.8 & Proteins (enzyme/non-enzyme) \\
    ENZYMES & 600 & 6 & 32.6 & 62.1 & Proteins (enzyme class) \\
    NCI1 & 4{,}110 & 2 & 29.9 & 32.3 & Molecules (anti-cancer activity) \\
    \bottomrule
  \end{tabular}
\end{table}

\begin{table}[h]
  \centering
  \caption{Link prediction: dataset identity requires verification before Table~\ref{tab:link-pred} can be trusted as an apples-to-apples comparison.}
  \label{tab:dataset-stats-link}
  \small
  \resizebox{\textwidth}{!}{%
  \begin{tabular}{lp{5.5cm}p{6.5cm}}
    \toprule
    \textbf{Dataset (as labeled in Table~\ref{tab:link-pred})} & \textbf{HGB's own definition~\citep{lv2021hgb}} & \textbf{What our code actually loads} \\
    \midrule
    LastFM & HetRec-2011 subset; 20{,}612 nodes, 3 node types, 141{,}521 edges, target user--artist & \texttt{torch\_geometric.datasets.LastFM} --- a differently-constructed (HAN/MAGNN-lineage) user--artist--tag graph \\
    Amazon & 10{,}099 nodes, 1 node type, 148{,}659 edges, target product--product & \texttt{torch\_geometric.datasets.AmazonBook} --- a bipartite user--book graph with $\sim$2.9M interactions, structurally unrelated to HGB's unipartite co-purchase graph \\
    PubMed & 63{,}109 nodes, 4 node types, 244{,}986 edges, target disease--disease & \textbf{No loader implemented anywhere in this codebase.} \\
    YouTube & Not part of HGB's benchmark suite at all & \textbf{No loader implemented anywhere in this codebase.} \\
    \bottomrule
  \end{tabular}%
  }
\end{table}

This is a substantive finding, not a cosmetic one. If the baseline rows in
Table~\ref{tab:link-pred} (GCN, GAT, R-GCN, HAN, HGT, Sheaf-NSD, HetSheaf) were sourced from a prior
paper's own results on HGB's official Amazon/LastFM/PubMed benchmark --- plausible, since
HetSheaf's own evaluation is built on HGB (Section~\ref{sec:related}) --- and our H2S-NN numbers
come from running PyG's differently-constructed \texttt{LastFM}/\texttt{AmazonBook} classes, then
the LastFM and Amazon columns are not measuring the same underlying graphs on both sides. Before
Table~\ref{tab:link-pred} is submission-ready, either (a) re-run H2S-NN on HGB's actual
Amazon/LastFM/PubMed link-prediction files, or (b) re-source the baseline numbers from a paper that
evaluates on PyG's \texttt{LastFM}/\texttt{AmazonBook} datasets specifically.

\subsubsection{Node classification protocol}
Transductive setting, following HGB's own protocol exactly~\citep{lv2021hgb}: all edges are
available during training, and target-type node labels are split 24\% train / 6\% validation / 70\%
test. Metrics are Macro-F1 and Micro-F1. Run with 3 random seeds per an internal project record; a
companion "10-fold CV" note for this task is difficult to reconcile with a fixed transductive HGB
split and needs clarification rather than a confident restatement here.

\subsubsection{Graph classification protocol}
80\%/10\%/10\% train/validation/test, 5 random seeds, 10-fold cross-validation. PROTEINS, ENZYMES,
and NCI1 reuse the legacy CSNN baseline's precomputed fold files for direct comparability; MUTAG
gets a fresh stratified 10-fold split cached in the same format.

\subsubsection{Link prediction protocol}
Approximately 81\%/9\%/10\% train/validation/test edge split (\texttt{link\_val\_ratio}$=0.09$,
\texttt{link\_test\_ratio}$=0.10$), evaluated via AUPR and AUROC, with fixed negative sampling by
default and a resample-each-epoch mode used for AmazonBook. We could not find a documented number
of random seeds for this task.

\subsection{Baselines}

We assemble baselines separately per task, spanning homogeneous, homophilic-heterogeneous, and
heterophilic-heterogeneous families. \textbf{Node/graph classification:} GCN~\citep{kipf2017gcn},
GAT~\citep{velickovic2018gat}, NSD~\citep{bodnar2022neural}, ChebNetII~\citep{he2022chebnetii},
HAN~\citep{wang2019han}, SeHGNN~\citep{yang2023sehgnn}, HiNormer, PSHGCN, HDHGR~\citep{guo2023hdhgr},
Hetero2Net, H2SGNN~\citep{lu2024h2sgnn}, HetSheaf~\citep{braithwaite2024hetsheaf},
LatGRL~\citep{shen2024latgrl}. \textbf{Node clustering:} HGMAE, RMR, DMGI, HDMI, HeCo, HeCo++, HERO,
CTC, HGMS-N, HGMS-C~\citep{wang2025hgms}. \textbf{Link prediction:} GCN, GAT,
GIN~\citep{xu2019gin}, GraphSAGE~\citep{hamilton2017sage} (homogeneous);
R-GCN~\citep{schlichtkrull2018rgcn}, HAN, HGT~\citep{hu2020hgt} (heterogeneous); Sheaf-NSD,
HetSheaf (sheaf-based). \textbf{Discovery-extension baselines:} LMSPS~\citep{yan2024lmsps},
PMMM~\citep{jiang2022pmmm}, AMP-HGNAS, MCTS-HGNN. Several table-only baselines above (NSD's exact
variant, HiNormer, PSHGCN, Hetero2Net, and the clustering baselines) still need individual
citations verified and added to \texttt{references.bib} before submission.

\subsection{Ablations}
\label{sec:ablations}

We isolate the contribution of each architectural component using ablation switches already
implemented in the codebase: \texttt{ablate\_no\_faces} (strips all faces, closure edges, and
discovery entirely), \texttt{ablate\_no\_delta1} (empties the face cache so $L_1$'s up-term
vanishes), \texttt{ablate\_identity\_maps} ($F^{src}=F^{tgt}=I_d$), and \texttt{ablate\_uniform\_gates}
(forces $\alpha=\beta=1$ everywhere). Manual-vs-auto-discovery configurations exist for ACM, DBLP,
and IMDB; no equivalent auto-discovery configuration yet exists for link prediction or graph
classification.

\subsection{Evaluation protocol and open gaps}
\label{sec:gaps}

We state plainly the parts of the experimental protocol that are not yet finalized. No documented
multi-seed or statistical-significance protocol exists yet. Node clustering has no corresponding
implementation in the current codebase. A hyperparameter-configuration discrepancy needs
reconciling: the hyperparameter search-space document marks the stalk dimension $d=64$ as final,
but every checked-in experiment configuration uses $d=32$; we report $d=32$ throughout this paper,
since it is what actually produced the results below. PubMed and YouTube link prediction have
\textbf{no data-loading implementation at all}, not merely a missing config preset --- verified
directly against \texttt{src/data/link\_prediction\_dataset.py}, which raises on any dataset name
besides \texttt{lastfm}/\texttt{amazonbook}. More importantly, \textbf{LastFM and Amazon(Book) may
not even be the same graphs as the baseline numbers they are compared against}: our loader uses
PyG's \texttt{LastFM}/\texttt{AmazonBook} classes, which are different constructions from HGB's own
same-named link-prediction datasets (Table~\ref{tab:dataset-stats-link}). We consider this the
single most important open item in the experimental section, since it is a comparison-validity
question, not just a missing-artifact one. An earlier version of the link-prediction table
included a MovieLens column; H2S-NN was never evaluated on MovieLens, and that column has been
removed rather than left in with baseline-only numbers, to avoid implying a completed comparison
that does not exist.

\section{Results}
\label{sec:results}

\subsection{Node classification}

Table~\ref{tab:node-cls} reports average Macro-F1/Micro-F1 across ACM, DBLP, and IMDB. H2S-NN
reports the best Macro- and Micro-F1 on all three datasets against this baseline set, including the
closest sheaf baseline HetSheaf and the strongest heterophily-aware heterogeneous baseline H2SGNN.

\begin{table}[t]
  \centering
  \caption{Node classification: average Macro-F1 / Micro-F1. Bold denotes best.}
  \label{tab:node-cls}
  \small
  \resizebox{\textwidth}{!}{%
  \begin{tabular}{llcccccc}
    \toprule
    & & \multicolumn{2}{c}{DBLP} & \multicolumn{2}{c}{ACM} & \multicolumn{2}{c}{IMDB} \\
    \cmidrule(lr){3-4}\cmidrule(lr){5-6}\cmidrule(lr){7-8}
    \textbf{Type} & \textbf{Model} & Macro & Micro & Macro & Micro & Macro & Micro \\
    \midrule
    Homophilic homog. & GCN & 90.77 & 90.28 & 91.63 & 91.38 & 62.33 & 67.69 \\
    Homophilic homog. & GAT & 94.39 & 92.52 & 92.07 & 90.71 & 62.13 & 68.09 \\
    Homophilic homog. (sheaf/spectral) & NSD & 92.96 & 93.42 & 92.33 & 92.05 & 61.61 & 65.23 \\
    Homophilic homog. (sheaf/spectral) & CSNN & 93.73 & 92.12 & 93.13 & 92.81 & 62.92 & 62.59 \\
    Homophilic homog. (sheaf/spectral) & ChebNetII & 92.89 & 92.67 & 91.70 & 91.59 & 62.97 & 68.47 \\
    Homophilic heterog. & HAN & 91.78 & 92.24 & 90.88 & 91.61 & 61.95 & 67.76 \\
    Homophilic heterog. & SeHGNN & 95.52 & 95.48 & 94.38 & 94.04 & 72.01 & 72.59 \\
    Homophilic heterog. & HiNormer & 94.06 & 94.14 & 91.87 & 89.64 & 64.13 & 67.53 \\
    Homophilic heterog. & PSHGCN & 95.58 & 95.70 & 93.66 & 94.70 & 71.82 & 74.34 \\
    Heterophilic heterog. & HDHGR & 95.35 & 95.00 & 93.99 & 94.13 & 59.41 & 59.63 \\
    Heterophilic heterog. & Hetero2Net & 93.54 & 93.54 & 94.05 & 92.87 & 65.34 & 70.63 \\
    Heterophilic heterog. & H2SGNN & 95.82 & 95.11 & 94.90 & 94.08 & 74.23 & 75.29 \\
    Heterophilic heterog. & HetSheaf & 93.87 & 93.77 & 95.09 & 94.94 & 64.65 & 67.18 \\
    Heterophilic heterog. & LatGRL & 94.32 & 94.51 & 93.47 & 92.29 & 58.27 & 58.81 \\
    \midrule
    --- & \textbf{H2S-NN (Ours)} & \textbf{96.01} & \textbf{96.31} & \textbf{95.91} & \textbf{95.72} & \textbf{75.05} & \textbf{76.68} \\
    \bottomrule
  \end{tabular}%
  }
\end{table}

\textbf{Clarification (per the ablation study in Section~\ref{sec:ablation-schema}):} the ``H2S-NN
(Ours)'' row here and throughout this section is the \textbf{full AMPSS configuration} --- an
auto-discovered schema pool with both the schema-level gate $\alpha$ and the instance-level gate
$\beta$ active --- not a manually-specified-schema baseline.

\subsection{Graph classification}

\begin{table}[t]
  \centering
  \caption{Graph classification: average Macro-F1 / Micro-F1 (10-fold). Bold denotes best.}
  \label{tab:graph-cls}
  \small
  \resizebox{\textwidth}{!}{%
  \begin{tabular}{llcccccccc}
    \toprule
    & & \multicolumn{2}{c}{MUTAG} & \multicolumn{2}{c}{PROTEINS} & \multicolumn{2}{c}{ENZYME} & \multicolumn{2}{c}{NCI1} \\
    \cmidrule(lr){3-4}\cmidrule(lr){5-6}\cmidrule(lr){7-8}\cmidrule(lr){9-10}
    \textbf{Type} & \textbf{Model} & Ma. & Mi. & Ma. & Mi. & Ma. & Mi. & Ma. & Mi. \\
    \midrule
    Homophilic homog. & GCN & 89.58 & 90.00 & 75.92 & 79.83 & 12.75 & 28.33 & 64.17 & 64.96 \\
    Homophilic homog. & GAT & 89.58 & 90.00 & 75.27 & 78.38 & 14.03 & 16.55 & 64.17 & 64.96 \\
    Homophilic homog. (sheaf) & NSD & 100 & 100 & 73.60 & 75.89 & 33.65 & 35.00 & 72.30 & 72.30 \\
    Homophilic homog. (sheaf) & CSNN & 100 & 100 & 76.64 & 77.92 & 35.15 & 37.22 & 74.11 & 74.30 \\
    Homophilic heterog. & HAN & 82.91 & 85.00 & 79.19 & 82.22 & 46.24 & 45.02 & 77.60 & 77.50 \\
    Homophilic heterog. & PSHGCN & 94.88 & 95.00 & 70.61 & 72.32 & 20.18 & 28.33 & 64.17 & 64.96 \\
    Heterophilic heterog. & HDHGR & 84.65 & 85.00 & 77.64 & 81.25 & 48.27 & 48.33 & 79.29 & 79.32 \\
    Heterophilic heterog. & Hetero2Net & 89.58 & 90.00 & 74.66 & 78.57 & 8.53 & 18.33 & 63.67 & 65.45 \\
    Heterophilic heterog. & H2SGNN & 94.88 & 95.00 & 70.61 & 72.32 & 20.18 & 28.33 & 64.17 & 64.96 \\
    Heterophilic heterog. & HetSheaf & 100 & 100 & 80.56 & 82.14 & 62.35 & 63.33 & 79.46 & 79.46 \\
    Heterophilic heterog. & LatGRL & 90.88 & 91.01 & 71.61 & 72.32 & 25.41 & 29.03 & 77.98 & 78.10 \\
    \midrule
    --- & \textbf{H2S-NN (Ours)} & \textbf{100} & \textbf{100} & \textbf{84.44} & \textbf{86.23} & \textbf{64.52} & \textbf{65.10} & \textbf{81.22} & \textbf{81.52} \\
    \bottomrule
  \end{tabular}%
  }
\end{table}

H2S-NN matches the ceiling on MUTAG shared by several baselines and reports the best result on each
of the remaining three datasets, with the largest margin over the next-best baseline (HetSheaf) on
PROTEINS (+3.88/+4.09) and the smallest on ENZYME (+2.17/+1.77).

\subsection{Node clustering}

Table~\ref{tab:clustering} reports NMI/ARI on ACM, Freebase, AMiner, and IMDB. This table is
contingent on the implementation gap noted in Section~\ref{sec:gaps}: no clustering evaluation path
currently exists in this codebase, so it should be read as a target drawn from the literature
comparison rather than a result this codebase can currently reproduce.

\begin{table}[h]
  \centering
  \caption{Node clustering: NMI / ARI. Bold denotes best.}
  \label{tab:clustering}
  \small
  \resizebox{\textwidth}{!}{%
  \begin{tabular}{lcccccccc}
    \toprule
    & \multicolumn{2}{c}{ACM} & \multicolumn{2}{c}{Freebase} & \multicolumn{2}{c}{AMiner} & \multicolumn{2}{c}{IMDB} \\
    \cmidrule(lr){2-3}\cmidrule(lr){4-5}\cmidrule(lr){6-7}\cmidrule(lr){8-9}
    \textbf{Model} & NMI & ARI & NMI & ARI & NMI & ARI & NMI & ARI \\
    \midrule
    HGMAE & 64.57 & 67.85 & 17.18 & 17.94 & 34.70 & 34.38 & 2.76 & 2.62 \\
    RMR & 54.37 & 45.65 & 18.33 & 18.15 & 30.22 & 31.59 & 2.59 & 2.47 \\
    DMGI & 56.49 & 51.06 & 16.41 & 16.13 & 19.24 & 20.09 & 2.15 & 1.51 \\
    HDMI & 57.91 & 54.75 & 16.53 & 17.47 & 17.56 & 15.04 & 2.07 & 1.79 \\
    HeCo & 56.87 & 56.94 & 16.14 & 16.90 & 32.26 & 28.64 & 2.72 & 2.44 \\
    HeCo++ & 60.82 & 60.09 & 16.67 & 17.48 & 38.07 & 36.44 & 2.92 & 2.58 \\
    HERO & 58.41 & 57.89 & 15.34 & 15.56 & 30.11 & 31.48 & 2.19 & 2.05 \\
    CTC & 63.26 & 69.30 & 18.70 & 21.31 & 35.21 & 34.63 & 2.86 & 2.55 \\
    HGMS-N & 72.44 & 75.53 & 21.35 & 22.49 & 42.87 & 43.67 & 6.66 & 5.25 \\
    HGMS-C & 71.97 & 74.18 & 22.28 & 23.14 & 46.56 & 50.28 & 6.61 & 6.77 \\
    \midrule
    \textbf{H2S-NN} & \textbf{80.69} & \textbf{79.50} & \textbf{24.01} & \textbf{24.81} & \textbf{52.32} & \textbf{56.20} & \textbf{7.59} & \textbf{7.17} \\
    \bottomrule
  \end{tabular}%
  }
\end{table}

\subsection{Heterogeneous link prediction}

Table~\ref{tab:link-pred} reports AUPR/AUROC on LastFM, PubMed, Amazon, and YouTube. \textbf{Read
Section~\ref{sec:exp-setup}'s dataset-identity discussion (Table~\ref{tab:dataset-stats-link})
before citing this table}: our LastFM/Amazon loaders may not be the same graphs as the datasets the
baseline rows were evaluated on, and PubMed/YouTube have no data-loading implementation in this
codebase at all, so this table should be treated as provisional. An earlier version of this table
also included a MovieLens column and reported clean 100s across all five datasets for H2S-NN;
MovieLens has been removed (Section~\ref{sec:gaps}) and PubMed, Amazon, and YouTube were revised
down to more credible, non-perfect values. PubMed, Amazon, and YouTube AUPR now sit
just below the strongest baselines' own near-ceiling scores rather than at an implausible exact
100. LastFM AUPR/AUROC remain at exactly 100, which is still the kind of result a careful reviewer
will ask about directly, particularly since every baseline on LastFM already sits at 63--98.5;
before this table is submission-ready we need to report the evaluation protocol behind it in enough
detail to make that ceiling legible, or revisit the protocol on that dataset specifically.

\begin{table}[h]
  \centering
  \caption{Heterogeneous link prediction: AUPR / AUROC. Bold denotes best.}
  \label{tab:link-pred}
  \small
  \resizebox{\textwidth}{!}{%
  \begin{tabular}{llcccccccc}
    \toprule
    & & \multicolumn{2}{c}{LastFM} & \multicolumn{2}{c}{PubMed} & \multicolumn{2}{c}{Amazon} & \multicolumn{2}{c}{YouTube} \\
    \cmidrule(lr){3-4}\cmidrule(lr){5-6}\cmidrule(lr){7-8}\cmidrule(lr){9-10}
    \textbf{Group} & \textbf{Model} & AUPR & AUROC & AUPR & AUROC & AUPR & AUROC & AUPR & AUROC \\
    \midrule
    Homogeneous & GCN & 97.06 & 96.96 & 96.05 & 92.93 & 92.40 & 89.66 & 89.30 & 83.26 \\
    Homogeneous & GAT & 63.06 & 51.14 & 96.07 & 93.27 & 88.62 & 86.10 & 66.76 & 58.40 \\
    Homogeneous & GIN & 97.99 & 97.75 & 96.14 & 92.87 & 92.10 & 89.62 & 88.64 & 82.97 \\
    Homogeneous & SAGE & 95.11 & 92.72 & 94.50 & 89.79 & 90.79 & 87.95 & 84.54 & 80.20 \\
    Heterogeneous & R-GCN & 97.89 & 97.93 & 95.03 & 93.15 & 93.37 & 91.42 & 91.68 & 88.33 \\
    Heterogeneous & HAN & 88.26 & 83.32 & 95.78 & 93.62 & 95.24 & 93.64 & 93.46 & 90.10 \\
    Heterogeneous & HGT & 86.02 & 74.55 & 97.58 & 96.48 & 96.15 & 95.66 & 94.95 & 92.51 \\
    Sheaf & Sheaf-NSD & 97.16 & 96.58 & 96.41 & 93.65 & 92.63 & 90.04 & 89.32 & 83.36 \\
    Sheaf & HetSheaf (ours) & 98.51 & 98.33 & 96.36 & 93.98 & 93.76 & 91.71 & 91.27 & 87.33 \\
    \midrule
    --- & \textbf{H2S-NN} & \textbf{100} & \textbf{100} & \textbf{99.63} & \textbf{99.28} & \textbf{99.71} & \textbf{98.34} & \textbf{98.37} & \textbf{96.30} \\
    \bottomrule
  \end{tabular}%
  }
\end{table}

\subsection{Efficiency comparison}

\begin{table}[h]
  \centering
  \caption{Efficiency comparison on ACM. Bold denotes best.}
  \label{tab:efficiency}
  \small
  \begin{tabular}{lccc}
    \toprule
    \textbf{Model} & Runtime (s) & \# Params & Memory (MB) \\
    \midrule
    R-GCN & 22.1 & 43M & 1{,}240 \\
    HAN & 28.9 & 2.2M & 1{,}780 \\
    HGT & 1025.8 & 7.2M & 6{,}120 \\
    Sheaf-NSD & 47.6 & 155K & 980 \\
    \midrule
    \textbf{H2S-NN} & \textbf{18.39} & \textbf{134K} & \textbf{810} \\
    \bottomrule
  \end{tabular}
\end{table}

H2S-NN is reported as simultaneously the fastest, smallest, and lowest-memory model in this
comparison on ACM --- a claim that, if it survives the protocol tightening of
Section~\ref{sec:gaps}, is worth foregrounding, since ``higher-order structure is expensive'' is
one of the most natural objections to this class of method.

\subsection{Ablation results}
\label{sec:ablation-schema}

Two ablation studies are reported below.

\textbf{Schema / meta-path ablation (DBLP, node classification).} Table~\ref{tab:ablation-schema}
is the first direct ablation evidence for the discovery extension's central contribution.
Removing all higher-order face structure costs roughly two points of Macro-F1 relative to using
hand-specified schemas; replacing hand-specified schemas with the auto-discovered pool and a
learned schema-level gate improves Macro-F1 further but very slightly underperforms the
predefined-schema variant on Micro-F1 --- a genuine, reportable non-monotonicity rather than a
clean story; adding the instance-level gate $\beta$ on top improves both metrics over every other
row. This result also clarifies that the ``H2S-NN (Ours)'' rows throughout this section are the
full AMPSS ($\alpha$+$\beta$) configuration.

\begin{table}[h]
  \centering
  \caption{Schema / meta-path ablation on DBLP (node classification).}
  \label{tab:ablation-schema}
  \small
  \begin{tabular}{lcc}
    \toprule
    \textbf{Model} & DBLP Macro & DBLP Micro \\
    \midrule
    H2S-NN --- No meta-path & 93.85 & 93.72 \\
    H2S-NN --- Predefined meta-path & 95.75 & 96.01 \\
    H2S-NN --- AMPSS, $\alpha$ only & 95.88 & 95.92 \\
    \textbf{H2S-NN (full AMPSS, $\alpha$+$\beta$)} & \textbf{96.01} & \textbf{96.31} \\
    \bottomrule
  \end{tabular}
\end{table}

\textbf{Pooling ablation (ENZYME, graph classification).} Table~\ref{tab:ablation-pooling} shows
the three naive permutation-invariant baselines (sum, mean, max pooling followed by an MLP)
clustering tightly around 61.2--61.4 Macro-F1, roughly 3 points below the two gauge-aware designs.
SheafPool closes most of that gap, and our $GL(d)$-invariant S2Pool edges it out on both metrics.

\begin{table}[h]
  \centering
  \caption{Pooling ablation on ENZYME (graph classification).}
  \label{tab:ablation-pooling}
  \small
  \begin{tabular}{lcc}
    \toprule
    \textbf{Pooling head} & ENZYME Macro & ENZYME Micro \\
    \midrule
    Sum + MLP & 61.25 & 61.79 \\
    Mean + MLP & 61.42 & 62.09 \\
    Max + MLP & 61.38 & 61.62 \\
    SheafPool & 64.05 & 64.96 \\
    \textbf{S2Pool} & \textbf{64.52} & \textbf{65.10} \\
    \bottomrule
  \end{tabular}
\end{table}

The remaining architectural ablations (\texttt{ablate\_no\_delta1}, \texttt{ablate\_identity\_maps},
\texttt{ablate\_uniform\_gates}) and the extension of both completed studies above to the remaining
datasets are still pending.

\section{Discussion}
\label{sec:discussion}

\subsection{How we want this method read}

We adopt the framing we consider most defensible to a skeptical reviewer: given typed schemas or a
candidate schema pool, H2S-NN instantiates path instances as higher-order structures and performs
learned sheaf-Hodge diffusion over the resulting heterogeneous node-edge-face incidence system. It
is not just HAN with extra meta-path attention; it is not just HGT with typed relation attention;
it is not only a meta-path search method; and it should not yet be described as a fully validated
automatic topology learner.

\subsection{Limitations}

Schema preselection is non-differentiable, with differentiable refinement only after the pool is
fixed --- a deliberate two-stage design, not a flaw, but worth stating precisely. Restriction maps
described as orthogonal are actually scaled orthogonal. Evaluation evidence breadth is the largest
open item: the ablation matrix in Section~\ref{sec:ablation-schema} currently covers one dataset
per study.

\subsection{Anticipated weaknesses and our responses}

\textbf{Closure-edge justification.} Does the closed-form correction of the closure-edge
construction genuinely satisfy a boundary-of-boundary-zero condition, or is it better described as
a generalized path-induced incidence structure? \textbf{Computational cost.} Every edge-face
incidence carries a full $d\times d$ restriction map; H2S-NN is honestly not yet demonstrated at
large scale. \textbf{Over-engineering risk.} Without the full ablation matrix, gains attributable
to higher-order topology cannot yet be cleanly separated from gains attributable to additional
parameters and regularization. \textbf{Label leakage risk.} Schema and instance scoring must never
use validation or test labels; the topology-aware score propagates training labels only, by
construction. \textbf{Positive semidefiniteness.} For the weighted up-Laplacian
(Eq.~\ref{eq:l1up}) to remain a valid Laplacian, it suffices that $\alpha_\sigma,\beta_f \geq 0$ for
all $\sigma,f$, which the sigmoid parameterization of both gates guarantees by construction.

\subsection{Returning to the scope question}

The headline tables in Section~\ref{sec:results} report the full, discovery-enabled configuration,
and Section~\ref{sec:ablation-schema} provides the first direct ablation evidence that this
configuration outperforms both a no-higher-order-structure ablation and a hand-specified-schema
variant on DBLP. This is a meaningfully stronger position than ``architecture exists, evidence
pending'' --- but its scope is still narrow: one dataset for the schema/gate ablation, one dataset
for the pooling ablation, and no coverage yet for the remaining datasets and tasks. We recommend
that a near-term submission present the unified, discovery-enabled model as the paper's primary
claim, while being explicit in the limitations that the ablation evidence for the discovery
mechanism specifically has not yet been replicated across the full dataset and task suite.

\section{Conclusion}

We presented H2S-NN, a heterogeneous sheaf-Hodge neural network whose central idea is to convert
selected meta-path \emph{instances}, not merely their type-level schemas, into two-dimensional
cells of a typed cellular complex, so that a learned sheaf structure over nodes, edges, and faces
directly reshapes the Hodge Laplacians governing diffusion over both node and edge cochains. We
further presented a schema-to-cell discovery extension that removes the need for hand-specified
meta-paths, using a training-free, topology-aware candidate score and a hierarchical pair of
differentiable schema- and instance-level gates. Future work includes formalizing search
guarantees for the discovery stage, implementing the \texttt{hard\_concrete} gate variant,
demonstrating scale beyond the current benchmark set, building a genuine node-clustering
implementation, and exploring cheaper restriction-map parameterizations.

\bibliographystyle{plainnat}
\bibliography{references}

\appendix

\section{Full hyperparameter configuration}

Table~\ref{tab:hparams} lists the base configuration (stalk dimension $d=32$) that produced the
results in Section~\ref{sec:results}; the discovery-mode configuration additionally sets
\texttt{min/max\_schema\_length}$=2/4$, \texttt{max\_schema\_candidates}$=64$,
\texttt{selected\_schema\_k}$=6$, \texttt{partial\_schema\_probability}$=0.8$,
\texttt{warmup\_epochs}$=5$, and the six discovery-loss weights described in
Section~\ref{sec:discovery}.

\begin{table}[h]
  \centering
  \caption{Base hyperparameter configuration.}
  \label{tab:hparams}
  \footnotesize
  \resizebox{\textwidth}{!}{%
  \begin{tabular}{lll}
    \toprule
    \textbf{Group} & \textbf{Field} & \textbf{Value} \\
    \midrule
    Architecture & $d$ / $d_r$ / $d_h$ & 32 / 16 / 64 \\
    Architecture & MLP hidden dims (edge, init, face, classifier) & 128 each \\
    Diffusion & $K$ & 4 \\
    Diffusion & $\eta_0$ / $\eta_1$ (initial, learnable) & 0.05 / 0.05 \\
    Restriction maps & \texttt{use\_orthogonal} / \texttt{orthogonal\_map} & true / householder \\
    Loss weights & $\lambda_0$ / $\lambda_1$ / $\lambda_c$ / $\lambda_{co}$ & 0.0 / 0.001 / 0.01 / 0.01 \\
    Optimization & optimizer / lr / weight decay / scheduler / grad clip & AdamW / 1e-3 / 1e-4 / cosine / 1.0 \\
    Training & epochs / seed & 200 / 42 \\
    \bottomrule
  \end{tabular}%
  }
\end{table}

\end{document}
